\section{Introduction}\label{sec:intro}

A state observer reconstructs the internal state of a dynamical system from its measured outputs. Kazantzis--Kravaris/Luenberger (KKL) observers \citep{kazantzis1998,andrieu2006existence} do this with a fixed linear filter, $\dot z = Az + By$ with $A$ Hurwitz, followed by the inverse of an injective map $T$ that sends the plant state into the filter's coordinates; existence of $T$ needs only mild observability conditions. Because $T$ is hard to compute in closed form, a line of work learns it and its inverse from simulated data \citep{ramos2020,janny2021,peralez2021,niazi2023,buisson2023,niazi2025kkl}. For plants driven by an exogenous input $u$, the theory says that the map must depend on the input history \citep{bernard2019luenberger}. It does not say where that dependence should be placed in a learned observer: in the latent dynamics, at the decoder's input, in its features, or in its weights. Nor does it say when such an observer is worth using, given that a practitioner with a model of the plant can run an extended Kalman filter (EKF) instead.

\paragraph{Why KKL observers.} Three properties set a KKL observer apart from a Kalman filter. Its latent filter does not depend on the estimate (for the decoder-side variants of Section~\ref{sec:framework}), so convergence is global and a poor initial guess cannot make it diverge, whereas the EKF's stability is local \citep{reif1999stochastic}; at deployment it needs no model and no Jacobians, only a linear filter and a feedforward decoder, which in our evaluation harness costs less per step than the EKF (Table~\ref{tab:scorecard}); and the map $T$ exists under mild observability, backward distinguishability \citep{andrieu2006existence,bernard2019luenberger}, including strongly nonlinear and chaotic plants where EKF linearisation is poor. Whether these advantages translate into accuracy is exactly what C3 tests.
% TODO (results layer): no per-step cost macro exists in numbers.tex yet (the us/step column of
% tab3_scorecard_realistic is the only generated source); quote the KKL/EKF cost ratio here once a
% macro (e.g. \CostEkfOverKkl) is generated. Do not hand-type it.

This paper answers both questions empirically, with a pre-registered protocol, and states the limits of the answers as plainly as the answers. HyperKKL, our open testbed, is the common framework: one latent filter, one pair of base maps trained once per plant, one fading-memory encoder of the input history, and five places where the resulting context can act. Our thesis is that KKL theory requires an input-dependent observer map but leaves the \emph{location} of that dependence and the \emph{value} of learning it open, and that both can be settled by a controlled experiment.

\paragraph{Contributions and verdicts.} Every claim below is tied to a decision rule written before the runs (Table~\ref{tab:decisions} in Appendix~\ref{app:tables}) and applied mechanically; failed rules are reported as such.
\begin{enumerate}[leftmargin=1.6em,itemsep=1pt]
\item \textbf{Where the input should enter (C1).} With identical frozen base maps, identical context encoder and a matched conditioner budget, HyperKKL, a hypernetwork that writes rank-$r$ updates of the decoder's weights from the input-history context, beats latent injection and an additive output residual on \ConeWeightSpaceBeatsLatConcShape\ dev plants for unseen waveform shapes and \ConeWeightSpaceBeatsLatConcSeq\ for unseen input sequences, and on \CtestWeightSpaceBeatsLatConcShape\ and \CtestWeightSpaceBeatsLatConcSeq\ held-out plants run once with the frozen recipe. The pre-registered rule asking for one location to beat \emph{every} other \FAIL s: HyperKKL and a dense-subspace ablation of it are tied, and the rule reported as passed is the weaker one pre-registered for the held-out plants.
\item \textbf{When learning beats a model-based filter (C3).} With an exact model the filters win in every one of fifteen stress conditions (sensor faults, input faults, poor initialisation, plant drift; learned-wins regime: \CthreeFaultsRegimes), and a tuned EKF is \ConeLearnedOverEkfMedianSeq$\times$ more accurate (median over dev plants; range \ConeLearnedOverEkfRangeSeq$\times$). Once the filter's model is wrong, the learned observers win: at every mismatch level tested, 10\% (\CthreeMismatchTenSeqVsWorking\ plants, sequence OOD), 20\% (\CthreeMismatchTwentySeqVsWorking), 30\% (\CthreeMismatchThirtySeqVsWorking), 50\% (\CthreeMismatchFiftySeqVsWorking); with one physical term missing (\CthreePartialSeqVsWorking); and in a realistic deployment combining a partial model, 30\% parameter error and coloured heavy-tailed noise (\CthreeRealisticSeqVsWorking\ against the best filter that still beats a constant predictor). The pre-registered crossover hypothesis, placed between 10\% and 30\% mismatch, \FAIL s in the direction favourable to learning (first winning level 10\%). The learned observers themselves fail under gross outliers and 30\% post-training drift on an exact model. On real battery drive cycles (Section~\ref{sec:realdata}) a pre-registered test confirms both orderings against a tuned EKF on an identified circuit model (\RealHOneLearnedBetter\ of \RealNTestCycles\ cycles each), with limits. This is two findings: (a) learning from real-plant data beats a wrong model, which is robust; (b) \emph{which} learner wins is system-dependent: on the partial-model node a grey-box EKF (model plus learned residual) or a Koopman Kalman filter trained on the same labels beats the KKL observers on three of five decided plants (Section~\ref{sec:c3}); the other nodes' same-data comparison is still being merged.
\item \textbf{A testbed and a way of reporting (C4).} Twelve plants with partial-knowledge models, an input protocol in which out-of-distribution means a new waveform shape or event sequence inside fixed actuator bounds, eight filters including robust and oracle variants, a climatology failure line, a deployment scorecard, and an experiment tree whose nodes carry their decision rules.
\item \textbf{Supporting analysis (C2).} A fading-memory account of how much input history the decoder needs (forgetting floor, memory probe, window sweeps); all its pre-registered rules \FAIL, and a refined rule restoring a plant-dependent constant explains part of the failures post hoc.
\end{enumerate}

\paragraph{Requirements and limits, up front.} (i) Every learned observer is trained on full-state labels from the real plant, so C3 compares ``data from the true plant'' with ``a wrong model'', not ``no knowledge'' with ``a model''. (ii) With an exact model a tuned EKF is several times better on every plant, so the learned observers' value lies in the model-wrong regime. (iii) Trained on the nominal sensor, they are fragile to gross outliers and post-training drift, where robust filters are not. (iv) The theory's counterfactual-past statements do not apply to the two chaotic plants; the real-data result holds at the training temperature only, on one cell, against a quickly identified circuit model. (v) The held-out network plant is weakly observable; the held-out C3 replication is \OutcomeTESTCThree\ (\JobsTESTCThree\ jobs); the same-data comparison is post hoc and, where decided, goes against the KKL observers.

Section~\ref{sec:framework} states the framework and three results of the theory informally; Section~\ref{sec:design} gives the experimental design; Sections~\ref{sec:c1}--\ref{sec:scorecard} report C1, C3, C2 and the scorecard; Section~\ref{sec:related} positions the work; Sections~\ref{sec:limitations}--\ref{sec:conclusion} state limitations, broader impact and conclusions. The proofs, the full tables and the history of the study are in the appendix.
